\section{Marco teórico}

\subsection{\normalsize{Fundamentos de información cuántica}}

\subsubsection{Qubits, espacio de estados y superposición}

La unidad básica de información en computación cuántica es el qubit. A diferencia del bit clásico, cuyo valor solo puede ser 0 o 1, un qubit se describe mediante un vector de estado en un espacio de Hilbert generado por los estados base $\ket{0}$ y $\ket{1}$ \cite{Nielsen2010}. Un estado general puede escribirse como
\[
    \ket{\psi} = \alpha \ket{0} + \beta \ket{1},
\]
donde $\alpha$ y $\beta$ son amplitudes complejas que satisfacen la condición de normalización
\[
    |\alpha|^2 + |\beta|^2 = 1.
\]

Esta expresión muestra que el estado cuántico no queda determinado por un valor discreto, sino por una combinación lineal de estados base. Esta propiedad se conoce como superposición y constituye una de las características fundamentales de la información cuántica \cite{Nielsen2010}. En términos observables, al medir el sistema se obtiene uno de los estados base con probabilidades dadas por los módulos cuadrados de las amplitudes.

La formulación anterior puede generalizarse a registros de varios qubits. En ese caso, el espacio de estados total se obtiene mediante el producto tensorial de los espacios individuales, de modo que un sistema de $n$ qubits se describe en un espacio vectorial de dimensión $2^n$ \cite{Nielsen2010}. Su estado general se expresa como
\[
    \ket{\psi} = \sum_{i=0}^{2^n-1} \alpha_i \ket{i},
\]
donde $\ket{i}$ representa cada estado base de la base computacional y $\alpha_i \in \mathbb{C}$. Por ejemplo, para dos qubits:
\[
    \ket{\psi} = \alpha_{00}\ket{00} + \alpha_{01}\ket{01} + \alpha_{10}\ket{10} + \alpha_{11}\ket{11}.
\]

La superposición adquiere aquí una forma más amplia: el sistema completo puede describirse como una combinación lineal de todos los estados base disponibles. Así, a medida que aumenta el número de qubits, crece también el número de amplitudes necesarias para describir el sistema con exactitud. Esta estructura matemática define el modo en que la información cuántica se representa antes de cualquier implementación computacional.

\subsubsection{Entrelazamiento}

El entrelazamiento aparece cuando el estado conjunto de un sistema no puede expresarse como el producto tensorial de estados independientes para cada una de sus partes \cite{Nielsen2010}. En estos casos, la información del sistema deja de poder describirse completamente a partir de sus componentes por separado y pasa a estar contenida en el estado global.

Un ejemplo clásico es el estado de Bell
\[
    \ket{\Phi^+} = \frac{1}{\sqrt{2}}\left(\ket{00} + \ket{11}\right).
\]
Este estado no puede factorizarse como $\ket{\psi_1} \otimes \ket{\psi_2}$, por lo que sus componentes no admiten una descripción independiente completa. El entrelazamiento introduce correlaciones no clásicas entre las partes del sistema y constituye uno de los rasgos distintivos de la computación cuántica \cite{Nielsen2010}.

Desde una perspectiva operativa, esta propiedad resulta fundamental porque muchas transformaciones cuánticas generan, modifican o explotan correlaciones de este tipo. En consecuencia, el entrelazamiento no solo es una característica teórica del modelo, sino también una propiedad estructural de gran importancia en la evolución de los estados cuánticos.

\subsubsection{Interferencia}

La interferencia es la propiedad mediante la cual las amplitudes complejas de un estado cuántico pueden combinarse de forma constructiva o destructiva durante la evolución del sistema \cite{Nielsen2010}. A diferencia de un esquema puramente probabilístico, en un sistema cuántico no se suman directamente probabilidades, sino amplitudes, y solo después se obtienen probabilidades a partir de sus módulos cuadrados.

Esto implica que la fase relativa entre amplitudes desempeña un papel central. Dos contribuciones de magnitud similar pueden reforzarse entre sí o cancelarse parcialmente según su relación de fase. Por esta razón, la evolución cuántica no depende únicamente de cuánto ``peso'' tenga cada estado base, sino también de cómo se relacionan sus amplitudes complejas.

La interferencia explica cómo ciertas configuraciones del sistema pueden aumentar su probabilidad de observación mientras otras disminuyen, aun cuando todas formen parte de la misma superposición. Junto con la superposición y el entrelazamiento, esta propiedad define la dinámica fundamental con la que la información cuántica es transformada.

\subsection{\normalsize{Modelo de circuitos cuánticos}}

\subsubsection{Evolución unitaria y compuertas cuánticas}

El modelo de circuitos cuánticos representa un algoritmo como una secuencia finita de transformaciones unitarias aplicadas sobre un registro de qubits \cite{Nielsen2010}. Si el estado actual del sistema es $\ket{\psi}$ y se aplica una operación unitaria $U$, el nuevo estado queda dado por
\[
    \ket{\psi'} = U\ket{\psi}.
\]

Las compuertas cuánticas son, por tanto, operadores unitarios que transforman el estado preservando su norma. Entre las compuertas de un qubit más utilizadas se encuentran las compuertas de Pauli, la compuerta Hadamard y distintas compuertas de fase.

La compuerta Hadamard se expresa como
\[
    H = \frac{1}{\sqrt{2}}
    \begin{bmatrix}
        1 & 1 \\
        1 & -1
    \end{bmatrix}
\]
y actúa sobre los estados base de la siguiente manera:
\[
    H\ket{0} = \frac{\ket{0} + \ket{1}}{\sqrt{2}}, \qquad
    H\ket{1} = \frac{\ket{0} - \ket{1}}{\sqrt{2}}.
\]

Esta compuerta es especialmente relevante porque transforma estados base en superposiciones y constituye un bloque elemental en múltiples circuitos.

En términos generales, si un circuito aplica sucesivamente las compuertas $U_1, U_2, \ldots, U_k$, el estado final puede escribirse como
\[
    \ket{\psi_f} = U_k \cdots U_2 U_1 \ket{\psi_0}.
\]
Esta descripción resalta que el circuito puede interpretarse como una composición ordenada de transformaciones lineales sobre el vector de estado.

\subsubsection{Compuertas controladas y operaciones sobre varios qubits}

Además de las operaciones de un solo qubit, los circuitos cuánticos utilizan compuertas que actúan sobre más de un qubit de manera conjunta. Una de las más conocidas es la compuerta CNOT, que aplica una negación controlada sobre un qubit objetivo dependiendo del valor del qubit de control \cite{Nielsen2010}.

Las compuertas controladas permiten construir correlaciones entre qubits y son fundamentales para la generación de entrelazamiento. Por ejemplo, si se parte de $\ket{00}$, se aplica una compuerta Hadamard al primer qubit y después una CNOT controlada por ese mismo qubit, el resultado es
\[
    \frac{\ket{00} + \ket{11}}{\sqrt{2}},
\]
que corresponde a un estado entrelazado.

En el modelo de circuitos, estas compuertas amplían la capacidad expresiva del sistema al introducir dependencias entre subsistemas. Su efecto puede describirse algebraicamente mediante matrices unitarias de mayor dimensión, aunque en la práctica suelen implementarse mediante reglas de actualización local sobre el vector de estado.

\subsection{\normalsize{Del modelo cuántico a su representación clásica}}

\subsubsection{El vector de estado como representación computacional}

La forma más directa de representar un sistema cuántico en una computadora clásica consiste en almacenar explícitamente su vector de estado. En este enfoque, el sistema de $n$ qubits se materializa como una secuencia de $2^n$ amplitudes complejas:
\[
    \ket{\psi} = \sum_{i=0}^{2^n-1} \alpha_i \ket{i}
    \longrightarrow
    [\alpha_0,\alpha_1,\alpha_2,\ldots,\alpha_{2^n-1}].
\]

Cada elemento del arreglo representa una amplitud del sistema, mientras que su posición identifica el estado base asociado. Esta equivalencia convierte la formulación matemática del estado cuántico en una estructura de datos concreta sobre la cual pueden aplicarse operaciones numéricas.

El uso explícito del vector de estado caracteriza a la simulación de estado completo, también conocida como simulación tipo Schr\"odinger \cite{Jones2019}. En ella, el sistema se representa manteniendo todas sus amplitudes y actualizándolas conforme se aplican las compuertas del circuito. La principal ventaja de esta aproximación es su correspondencia directa con la formulación matemática del modelo cuántico.

Existen otros enfoques de simulación clásica, como los basados en redes tensoriales, diagramas de decisión o métodos de suma sobre caminos \cite{Vidal2003,Viamontes2003,Chen2018}. De manera general, estos métodos buscan representar la evolución cuántica mediante estructuras alternativas que explotan diferentes formas de organización algebraica o combinatoria. Sin embargo, la representación explícita por vector de estado ofrece el marco más directo para comprender cómo se traduce un sistema cuántico a memoria clásica y cómo se implementa su evolución sobre arreglos numéricos. Esta distinción también es importante porque no todas las estrategias para reducir memoria cambian la representación matemática del estado: algunas, como la que se estudia en este trabajo, mantienen el formalismo Schr\"odinger y actúan únicamente sobre la manera de almacenar los datos.

\subsubsection{Aplicación de compuertas sobre el arreglo de amplitudes}

Una vez que el estado cuántico se representa como un arreglo lineal, la aplicación de compuertas puede entenderse como una secuencia de operaciones numéricas sobre subconjuntos específicos de amplitudes.

En una compuerta de un solo qubit, la transformación no afecta una amplitud aislada, sino pares de amplitudes cuyos índices difieren en el bit asociado al qubit objetivo. En un sistema de tres qubits puede tomarse la convención $(q_2 q_1 q_0)$, donde $q_2$ es el bit más significativo y $q_0$ el menos significativo. Con esta convención, la Tabla~\ref{tab:pares-bits-compuerta} muestra cómo cambian los pares afectados según la posición del qubit objetivo.

\begin{table}[H]
    \centering
    \footnotesize
    \begin{tabular}{c c c}
        \hline
        \multicolumn{3}{c}{Objetivo $q_0$ (se mantienen fijos $q_2$ y $q_1$)} \\
        \hline
        Bits fijos $(q_2 q_1)$ & Estados relacionados & Par afectado \\
        \hline
        $00$ & $000 \leftrightarrow 001$ & $(\alpha_0,\alpha_1)$ \\
        $01$ & $010 \leftrightarrow 011$ & $(\alpha_2,\alpha_3)$ \\
        $10$ & $100 \leftrightarrow 101$ & $(\alpha_4,\alpha_5)$ \\
        $11$ & $110 \leftrightarrow 111$ & $(\alpha_6,\alpha_7)$ \\
        \hline
    \end{tabular}

    \vspace{0.7em}

    \begin{tabular}{c c c}
        \hline
        \multicolumn{3}{c}{Objetivo $q_2$ (se mantienen fijos $q_1$ y $q_0$)} \\
        \hline
        Bits fijos $(q_1 q_0)$ & Estados relacionados & Par afectado \\
        \hline
        $00$ & $000 \leftrightarrow 100$ & $(\alpha_0,\alpha_4)$ \\
        $01$ & $001 \leftrightarrow 101$ & $(\alpha_1,\alpha_5)$ \\
        $10$ & $010 \leftrightarrow 110$ & $(\alpha_2,\alpha_6)$ \\
        $11$ & $011 \leftrightarrow 111$ & $(\alpha_3,\alpha_7)$ \\
        \hline
    \end{tabular}
    \caption{Pares de amplitudes afectados por una compuerta de un solo qubit en un sistema de 3 qubits, mostrando qué bits permanecen fijos y cuál varía según el qubit objetivo.}
    \label{tab:pares-bits-compuerta}
\end{table}

Esto muestra que el patrón de actualización depende de la posición del qubit sobre el cual se aplica la compuerta. De forma análoga, una compuerta de dos qubits modifica grupos más grandes de amplitudes, pero siempre siguiendo una regla determinada por la estructura binaria de los índices.

Desde el punto de vista computacional, esta regularidad introduce patrones de acceso sobre el arreglo. Algunas operaciones actúan sobre posiciones contiguas y otras sobre posiciones separadas por saltos regulares, de modo que la simulación no solo depende del tamaño total del vector, sino también de cómo se distribuyen las lecturas y escrituras en memoria. En este contexto, las nociones de localidad espacial y localidad temporal ayudan a entender por qué ciertas estrategias de particionamiento, acceso selectivo y reutilización temporal pueden resultar convenientes.

Esta interpretación computacional del circuito es importante porque establece el vínculo entre álgebra lineal y acceso a memoria: simular la evolución cuántica equivale, en este marco, a recorrer, leer, combinar y reescribir posiciones concretas de un arreglo de amplitudes complejas.

\subsubsection{Consideraciones básicas de almacenamiento y acceso}

Si cada amplitud compleja ocupa 16 bytes, el consumo de memoria del vector de estado puede expresarse como
\[
    M(n)=2^n \times 16\ \text{bytes}.
\]
Esta relación resume el costo directo de almacenar explícitamente el sistema en una simulación de estado completo \cite{Jones2019,wu2019full}. 

Por tanto, la representación clásica del vector de estado involucra dos aspectos complementarios: cuánto espacio se necesita para almacenar las amplitudes y cómo se accede a ellas durante la evolución del circuito. Estos dos elementos son fundamentales para entender la gestión de memoria sobre arreglos numéricos de gran tamaño.

\subsection{\normalsize{Gestión de memoria y compresión de datos en arreglos numéricos}}

\subsubsection{Gestión de memoria en simulación científica}

La gestión de memoria en computación científica comprende las técnicas utilizadas para organizar, almacenar, recuperar y reutilizar estructuras numéricas de gran tamaño durante la ejecución de un programa. En este tipo de aplicaciones, el rendimiento no depende únicamente del costo aritmético de las operaciones, sino también del comportamiento de los datos en memoria.

Cuando una estructura domina el volumen de almacenamiento y concentra buena parte de los accesos, su organización interna se vuelve especialmente importante. En estos casos, aspectos como la granularidad del acceso, la reutilización temporal y la disposición lineal o segmentada de los datos pueden influir de manera significativa en el comportamiento global del sistema.

Desde esta perspectiva, una estructura como el vector de estado puede analizarse no solo como objeto matemático, sino también como carga principal de memoria dentro de una simulación numérica.

\subsubsection{Conceptos generales de compresión de datos}

La compresión de datos es el proceso mediante el cual una secuencia de información se representa utilizando menos espacio que en su forma original. La posibilidad de comprimir depende de la existencia de redundancia o regularidad aprovechable en los datos.

En términos generales, los métodos de compresión se clasifican en sin pérdida y con pérdida. La compresión sin pérdida permite recuperar exactamente los datos originales después de descomprimir. La compresión con pérdida, en cambio, admite una reconstrucción aproximada a cambio de mayores reducciones de tamaño.

En datos numéricos, la compresibilidad puede venir de repeticiones, regiones homogéneas, distribuciones estructuradas o relaciones locales entre valores cercanos. Sin embargo, los arreglos en punto flotante presentan una estructura binaria que suele ser menos intuitiva que la de datos textuales o discretos, por lo que su compresión requiere técnicas adecuadas a esa naturaleza.

\subsubsection{Compresión sin pérdida en datos de punto flotante}

La compresión sin pérdida de arreglos de punto flotante busca reducir el espacio de almacenamiento preservando exactamente cada valor original. Esto implica que, tras la descompresión, tanto la representación numérica como la secuencia binaria asociada a cada dato deben coincidir con las del arreglo inicial.

Aunque los datos en punto flotante pueden parecer poco compresibles a simple vista, en muchos contextos científicos presentan regularidades explotables. Estas regularidades no siempre se manifiestan como repeticiones exactas de valores, sino también como similitudes locales en exponentes, mantisas o agrupaciones de datos con cierta estructura. Compresores como FPC y FPZIP fueron diseñados precisamente para explotar este tipo de regularidades en arreglos de punto flotante sin alterar sus valores \cite{burtscher2009fpc,lindstrom2006fpzip}.

Por ello, la compresibilidad de un arreglo numérico depende tanto del contenido como de la forma en que este se organiza antes de aplicar el compresor. Esta observación es importante cuando se trabaja con arreglos densos de amplitudes complejas, cuya estructura puede variar según el estado representado.

\subsubsection{Compresión por bloques, descompresión selectiva y caché}

Una estrategia habitual para manejar arreglos grandes consiste en particionarlos en bloques y comprimir cada bloque de manera independiente. Esta organización permite trabajar de forma más granular sobre los datos y evita la necesidad de comprimir o descomprimir la estructura completa ante cada operación. En propuestas recientes de simulación cuántica con compresión, este tipo de organización se combina con mecanismos de acceso parcial y administración jerárquica de memoria para reducir presión sobre la RAM \cite{wu2019full,zhang2024bmqsim}.

La descompresión selectiva se basa en esa granularidad: solo se reconstruyen temporalmente los bloques necesarios para una operación determinada. De este modo, la compresión puede integrarse con el patrón de acceso del sistema y no únicamente con el almacenamiento final de los datos.

La granularidad del bloque introduce, sin embargo, un compromiso importante. Bloques grandes suelen capturar mejor la redundancia disponible y favorecer mejores tasas de compresión, pero incrementan el costo cuando solo una pequeña fracción de los datos es necesaria. Bloques pequeños permiten accesos más localizados, aunque pueden degradar el rendimiento del compresor y aumentar el peso relativo del metadato y de la administración.

Cuando ciertos bloques son reutilizados con frecuencia en intervalos cortos, resulta útil mantenerlos temporalmente en una caché. Una política frecuente para administrar esta caché es LRU (\emph{Least Recently Used}), que reemplaza primero los bloques menos recientemente utilizados. Esta política se apoya en la localidad temporal y busca reducir el costo de descompresiones repetidas.

Bloques, descompresión selectiva y caché conforman así un conjunto de conceptos que permiten pensar la compresión como parte activa de la gestión de memoria sobre arreglos numéricos de gran tamaño.

\subsubsection{Exactitud numérica e igualdad exacta}

En aplicaciones científicas, no siempre basta con obtener resultados aproximados. En muchos casos se requiere preservar exactamente los valores almacenados, de modo que cualquier transformación aplicada a los datos sea completamente reversible.

La compresión sin pérdida satisface este requisito mediante la propiedad de igualdad exacta, según la cual los datos recuperados después de descomprimir coinciden completamente con los originales. Este criterio es más fuerte que otras métricas numéricas basadas en error absoluto, error relativo o fidelidad, ya que no admite desviaciones, por pequeñas que sean. Expresado sobre un vector de amplitudes, este requisito puede escribirse como
\[
    \alpha_i^{(\text{orig})} = \alpha_i^{(\text{rec})} \quad \forall i.
\]

En arreglos de amplitudes complejas, esta propiedad significa que tanto la parte real como la parte imaginaria de cada elemento deben recuperarse sin alteración. Como medida auxiliar de discrepancia, puede emplearse el error absoluto máximo
\[
    e_{\text{abs}} = \max_i \left|\alpha_i^{(\text{orig})} - \alpha_i^{(\text{rec})}\right|,
\]
el cual debe ser nulo en una estrategia estrictamente sin pérdida. Desde un punto de vista metodológico, ello permite distinguir claramente entre optimización del almacenamiento y modificación del contenido numérico.
